% explainer-source-hash: 8e875e65cc04b8f0bc7857f6ff372989bf03e63de2b94b820d9e2552fda1341c
% explainer-source-commit: 623a91e76858f28487b4e9b8ffdc4494380f8889
% explainer-generated: 2026-07-16
% Regenerate with /explain-all (see WIRING.md). Do not edit this stamp by hand:
% it is how the toolchain knows whether this document still matches the code.
\documentclass[11pt,a4paper]{article}
\input{preamble}

\title{\vspace{-1.5cm}\textbf{Codeforces Mashup Generator}\\[0.3cm]
\large A Deep Dive, from zero prior knowledge}
\author{Explainer for Salman Adnan's portfolio}
\date{}

\begin{document}
\maketitle
\thispagestyle{empty}

\begin{keyidea}{The whole project in four sentences}
A study group wanted to run practice contests, but people kept being handed problems
they had already solved, which defeats the point of a contest. This tool asks the public
Codeforces website for its entire problem catalogue, then asks it for the full
submission history of every group member, and throws away every problem that any member
has already solved. It can also insist that at least one strong, trusted competitor has
solved a problem before it is considered worth setting. What survives is written to a
spreadsheet, grouped by difficulty.
\end{keyidea}

\tableofcontents
\newpage

\section{Orientation: what this thing actually is}

\subsection{The vocabulary you need first}

This project sits inside a hobby, competitive programming, that has its own dense
vocabulary. Nothing below is assumed.

\begin{jargon}{Competitive programming}
A sport where people solve small, self-contained algorithm puzzles against the clock.
You are given a problem statement, you write a program, you submit it, and an automated
judge either accepts it or rejects it.
\end{jargon}

\begin{jargon}{Codeforces}
One of the largest websites hosting this sport. It stores tens of thousands of problems
and a permanent record of every submission every user has ever made. It exposes that
record publicly through an API (defined below), which is the only reason this project can
exist.
\end{jargon}

\begin{jargon}{API (Application Programming Interface)}
A door into a website meant for programs rather than people. Instead of a page of styled
text for a human to read, you request a URL and get back structured data for a program to
read. Codeforces' API is public: no key, no password, no account needed.
\end{jargon}

\begin{jargon}{JSON}
The format that structured data comes back in. It is text made of nested lists and
labelled fields, which a program can turn directly into its own native lists and
dictionaries. A Codeforces problem arrives as JSON that looks, in spirit, like: contest
number 2236, letter A, name ``Games on the Train'', rating 800, tags greedy and math.
\end{jargon}

\begin{jargon}{Rating (of a problem)}
Codeforces labels each problem with a difficulty number, roughly 800 (easiest) to 3500
(hardest). This project's default range is 800 to 2000 in steps of 100. Rating is the
axis the whole tool is organised around: the output is ``give me N problems at each of
these difficulties''.
\end{jargon}

\begin{jargon}{Tags}
Codeforces labels each problem with the techniques needed to solve it: \code{greedy},
\code{dp} (dynamic programming), \code{bitmasks}, \code{graphs}, and so on. A problem
usually carries several at once.
\end{jargon}

\begin{jargon}{Mashup}
A Codeforces feature that lets anyone assemble a private practice contest out of existing
problems, rather than writing new ones. The tool's whole job is choosing which existing
problems go in. It does not create the mashup itself; a human takes the spreadsheet and
sets it up.
\end{jargon}

\begin{jargon}{Verdict, and \code{OK}}
The judge's ruling on a submission. \code{OK} means accepted, that is, actually solved.
Every other verdict (wrong answer, time limit exceeded, compile error) means the person
tried and failed. This distinction is load-bearing: the tool counts a problem as ``solved
by you'' only on an \code{OK} verdict, so problems you attempted and failed are still
fair game.
\end{jargon}

\subsection{The problem being solved}

The tool answers one question, per difficulty band:

\begin{quote}
\itshape Which Codeforces problems has nobody in my group solved, that a strong
competitor has solved?
\end{quote}

The two halves of that question do different jobs, and it is worth being precise about
which is which:

\begin{itemize}
\item The \textbf{exclude} half is the point of the tool. It enforces fairness. If any
group member has an accepted submission on a problem, that problem is out.
\item The \textbf{include} half is a quality heuristic, and only a heuristic. The
reasoning encoded in the default list is that if a well-known strong competitor has
bothered to solve a problem, it is probably a real problem worth setting rather than an
obscure or broken one. This is a proxy, not a measure of quality, and section~\ref{sec:honest}
is blunt about what it costs.
\end{itemize}

\begin{honest}{The include filter is expensive for what it buys}
The include filter's stated purpose is a rough quality signal. Codeforces already ships a
direct quality signal in the problemset response itself: \code{solvedCount}, the number of
people who have solved each problem, which arrives free with data the tool already
downloads. The tool instead downloads the complete lifetime submission history of 46
separate people to approximate the same idea. The README's own ``What I'd do differently''
section names this, which is to its credit. I measured the actual cost in
section~\ref{sec:perf}: it is worse than the README implies.
\end{honest}

\section{The shape of the repository}

The entire project is one Python file. Everything else is packaging, documentation, or, in
one case, dead weight.

\begin{figure}[H]
\centering
\begin{forest} dirtree
[codeforces-mashup-generator
  [generate\_mashup.py {\rmfamily\itshape\quad the entire program, 228 lines}]
  [requirements.txt {\rmfamily\itshape\quad three dependencies}]
  [all\_users.txt {\rmfamily\itshape\quad \textcolor{warnred}{dead: never read by anything}}]
  [.env.example {\rmfamily\itshape\quad states that no env vars exist}]
  [.gitignore]
  [LICENSE {\rmfamily\itshape\quad PolyForm Strict 1.0.0}]
  [README.md]
  [docs
    [screenshots
      [app.png]
    ]
    [explainers
      [deep-dive.tex {\rmfamily\itshape\quad this document}]
      [brief.tex]
    ]
  ]
]
\end{forest}
\caption{The tracked files. One script, no package structure, no test suite.}
\end{figure}

\begin{keyidea}{Why a single file is the right call here}
A beginner is often told to split code into modules. That advice has a threshold. This
program is a linear pipeline with five functions, no state that outlives a run, and no
second consumer of any of its parts. Splitting it across files would add navigation cost
and buy nothing. The absence of structure here is proportionate, not lazy.
\end{keyidea}

\subsection{The one file that should not be there}

\begin{honest}{\code{all\_users.txt} is dead data, and I can prove it is stale}
\code{all\_users.txt} sits at the project root holding two Python lists,
\code{solver\_handles} and \code{exclude\_handles}. It looks like configuration. It is
not. \code{generate\_mashup.py} performs \textbf{no file I/O whatsoever}: it has no
\code{open()} call, no file read, and no import of this file. The script's own
\code{DEFAULT\_INCLUDE\_USERS} and \code{DEFAULT\_EXCLUDE\_USERS} constants are what
actually run. Deleting \code{all\_users.txt} would change nothing.

I compared the two mechanically rather than by eye, and the result pins down its history:
\begin{itemize}
\item The include lists are \textbf{byte-for-byte identical}: 46 handles, same order.
\item The exclude lists are \textbf{not}. The file has 28 handles; the script has 34. The
file's 28 are a \textbf{strict subset} of the script's 34, missing exactly these six:
\code{salmanadnan2025}, \code{Aritro\_}, \code{Drakozs}, \code{chpp}, \code{Vyaduct},
\code{chiru200513}. Nothing is in the file that is not in the script.
\end{itemize}
A strict subset in one list and an exact match in the other is the signature of a snapshot
that was forked and then left behind while the live copy kept growing by six members.
\textbf{Confirmed:} the file is unread and its exclude list is a strict subset.
\textbf{Inferred:} that it is a pre-refactor leftover rather than something intended for
future use. The code cannot tell us intent. It should be deleted; a file that looks like
config but is not is a trap for the next reader, and this is the second document to have
to explain it.
\end{honest}

\section{Architecture}

\subsection{The pipeline}

The program is a straight line. Data enters from two API endpoints, gets reduced to sets,
gets filtered, gets truncated, and lands in a spreadsheet. There are no branches in the
overall shape, no loops back, and no persistence between runs.

\begin{figure}[H]
\centering
\begin{tikzpicture}[x=1cm, y=1cm]
  % lane L = the catalogue path, lane R = the submissions path
  \node[extbox] (cfp) at (0,0) {Codeforces API\\\code{problemset.problems}};
  \node[extbox] (cfs) at (7.4,0) {Codeforces API\\\code{user.status}};

  \node[box] (fetchp) at (0,-2.0) {\code{fetch\_problemset}\\\scriptsize one request, 60s timeout};
  \node[box] (fetchu) at (7.4,-2.0) {\code{fetch\_user\_status}\\\scriptsize paginated, per handle};

  \node[store, minimum height=12mm] (byrating) at (0,-4.2) {\code{by\_rating}\\\scriptsize rating $\to$ problems};
  \node[box] (collect) at (7.4,-4.2) {\code{collect\_solved}\\\scriptsize keep verdict OK only};

  \node[store] (sets) at (7.4,-7.3) {\code{include\_solved} / \code{exclude\_solved}\\\scriptsize sets of \code{(contestId, index)}};

  \node[dec] (gate) at (3.7,-9.7) {in include set?\\not in exclude set?\\tags match?};

  \node[box] (pick) at (3.7,-11.8) {\code{pick\_choices\_per\_rating}\\\scriptsize first N per rating};
  \node[box] (save) at (3.7,-13.4) {\code{save\_to\_excel}\\\scriptsize pandas $+$ openpyxl};
  \node[store] (xlsx) at (3.7,-15.0) {\file{.xlsx}};

  \draw[flow] (cfp) -- (fetchp);
  \draw[flow] (cfs) -- (fetchu);
  \draw[flow] (fetchp) -- (byrating);
  \draw[flow] (fetchu) -- (collect);
  \draw[flow] (collect) -- node[lbl, left=1mm, pos=0.5] {called once\\per handle list} (sets);
  \draw[flow] (byrating) |- (gate.west);
  \draw[dflow] (sets) |- node[lbl, pos=0.75] {membership\\tests} (gate.east);
  \draw[flow] (gate) -- node[lbl] {passes} (pick);
  \draw[flow] (pick) -- (save);
  \draw[flow] (save) -- (xlsx);
\end{tikzpicture}
\caption{The end-to-end pipeline, read top to bottom. The left lane fetches the catalogue,
the right lane fetches submission history and reduces it to sets of keys; they meet at the
filter. Dashed arrows are set-membership tests, solid arrows are data movement.}
\label{fig:arch}
\end{figure}

\subsection{The idea that makes it work}

\begin{keyidea}{One identity, normalised from two different envelopes}
This is the single design decision worth defending in an interview, and it is genuinely
the crux of the project.

Codeforces hands you the same logical thing, ``a problem'', in two \textbf{different
shapes} depending on which endpoint you ask.
\begin{itemize}
\item \code{problemset.problems} returns problem objects directly: \code{contestId} and
\code{index} sit at the top level.
\item \code{user.status} returns \emph{submission} objects, which wrap the problem one
level down under a \code{problem} key, alongside the verdict and timestamp.
\end{itemize}
You cannot compare those two objects. They are not the same shape, and even the parts
that overlap carry extra fields that differ per submission. So the code picks a
\textbf{canonical identity} that both can be reduced to: the tuple
\code{(contestId, index)}, for example \code{(2236, "A")}. Every submission collapses to
that tuple, every catalogue entry collapses to that tuple, and now the two sources are
directly comparable.

Once both sides speak the same identity, the entire filtering question becomes set
arithmetic: is this key in the include set, and is it absent from the exclude set. That is
why the code has no nested search anywhere.
\end{keyidea}

\begin{jargon}{Set, and why it is used here}
A set is a collection that answers ``is this item in you?'' in effectively constant time,
no matter how large it is, because it is indexed by a hash rather than scanned. A plain
list answers the same question by checking every element in turn. With roughly 11{,}000
problems tested against hundreds of thousands of submissions, the difference is between a
run that finishes and one that does not. The code builds \code{include\_solved} and
\code{exclude\_solved} as sets once, then does one O(1) lookup per problem.
\end{jargon}

\section{Walkthrough of every function}

The file has eight top-level definitions. I take them in dependency order rather than file
order.

\subsection{\code{parse\_args}: the entire user interface}

The tool has no configuration file and no interactive prompts. Every knob is a command
line flag, defined in one \code{argparse} parser.

\begin{jargon}{\code{argparse}}
Python's standard library for command line interfaces. You declare the flags you accept
and their defaults; it handles parsing, type conversion, validation, and generating the
\code{--help} text for free. The \code{ArgumentDefaultsHelpFormatter} used here also
prints each default into the help output, so \code{--help} is self-documenting.
\end{jargon}

\begin{table}[H]
\centering\small
\begin{tabular}{@{}llp{6.3cm}@{}}
\toprule
\textbf{Flag} & \textbf{Default} & \textbf{Effect} \\
\midrule
\code{--include-users} & 46 handles & Problem must be solved by at least one of these \\
\code{--exclude-users} & 34 handles & Problem must be solved by none of these \\
\code{--no-include-filter} & off & Skip the include stage entirely, and skip its download \\
\code{--no-exclude-filter} & off & Skip the exclude stage entirely \\
\code{--ratings} & 800..2000 by 100 & Which difficulty bands to fill \\
\code{--choices-per-rating} & 10 & How many options per band \\
\code{--include-tags} & (empty) & Tags to require, per \code{--tag-mode} \\
\code{--exclude-tags} & (empty) & Tags that disqualify a problem outright \\
\code{--tag-mode} & \code{all} & \code{off}, \code{some}, \code{all}, or \code{strict} \\
\code{--lang} & \code{en} & Statement language passed to the API \\
\code{--delay} & 0.5 & Seconds slept between API calls \\
\code{-o}, \code{--output} & \file{filtered\_problems.xlsx} & Output filename \\
\bottomrule
\end{tabular}
\caption{The complete CLI surface. Verified against the source, not the README.}
\end{table}

\begin{honest}{A default that contradicts itself, harmlessly}
\code{--tag-mode} defaults to \code{all}, but \code{--include-tags} defaults to empty, and
\code{passes\_tag\_filter} only consults the mode when the include-tags list is non-empty.
So the default mode is inert: it describes how to match a list that is not there. This is
not a bug, and it is arguably the friendliest default (the moment you pass tags, ``all''
is the intuitive reading), but it does mean the \code{--help} output advertises an active
default that does nothing on a bare run.
\end{honest}

\subsection{\code{fetch\_user\_status}: pagination, and knowing when to stop}

This is the most intricate function in the file, because it is the only one with a loop
whose exit condition is not obvious.

\begin{jargon}{Pagination}
An API that would otherwise return an enormous response instead returns it in slices. You
ask for a window (``give me 10{,}000 records starting at record 1''), then ask again for
the next window, until you have everything. Codeforces exposes this through \code{from}
and \code{count} parameters on \code{user.status}. The awkward part is always the same:
knowing when you have reached the end.
\end{jargon}

\begin{figure}[H]
\centering
\begin{tikzpicture}[x=1cm, y=1cm]
  % main spine, top to bottom
  \node[box] (start) at (0,0) {start: \code{start = 1}, \code{all\_subs = []}};
  \node[box] (req) at (0,-1.7) {GET \code{user.status?handle=H}\\\code{\&from=start\&count=10000}};
  \node[dec] (net) at (0,-3.6) {network or\\HTTP error?};
  \node[dec] (apiok) at (0,-5.6) {JSON \code{status}\\is \code{OK}?};
  \node[dec] (empty) at (0,-7.6) {page empty?};
  \node[box] (extend) at (0,-9.4) {append page to \code{all\_subs}\\then \code{sleep(delay)}};
  \node[dec] (short) at (0,-11.3) {page shorter\\than 10000?};
  \node[box, fill=softgrey] (ret) at (0,-13.3) {\textbf{return \code{all\_subs}}};

  % error boxes on the left, equal width so they share one return rail
  \node[box, draw=warnred, minimum width=34mm] (perr) at (-5.2,-3.6) {print error\\\scriptsize partial data kept};
  \node[box, draw=warnred, minimum width=34mm] (aerr) at (-5.2,-5.6) {print API comment\\\scriptsize partial data kept};

  % loop-back box on the right
  \node[box] (adv) at (4.6,-11.3) {\code{start += 10000}};

  \draw[flow] (start) -- (req);
  \draw[flow] (req) -- (net);
  \draw[flow] (net) -- node[lbl] {no} (apiok);
  \draw[flow] (apiok) -- node[lbl] {yes} (empty);
  \draw[flow] (empty) -- node[lbl] {no} (extend);
  \draw[flow] (extend) -- (short);
  \draw[flow] (short) -- node[lbl] {yes} (ret);

  \draw[flow] (net) -- node[lbl] {yes} (perr);
  \draw[flow] (apiok) -- node[lbl] {no} (aerr);

  % both error paths drop down a shared rail well clear of the boxes, into the return
  \draw[flow] (perr.west) -- ++(-0.9,0) |- (ret.west);
  \draw[flow] (aerr.west) -- ++(-0.9,0) |- (ret.west);

  % normal end-of-data exit, down the inner right lane
  \draw[flow] (empty.east) -- ++(1.5,0) |- (ret.east);

  % loop back up the outer right lane
  \draw[flow] (short) -- node[lbl, pos=0.45] {no} (adv);
  \draw[flow] (adv.east) -- ++(0.9,0) |- (req.east);
\end{tikzpicture}
\caption{\code{fetch\_user\_status} as a state machine. Note that all three red and grey
exits converge on the same return, and none of them signals failure to the caller. That
convergence is the root of the finding in section~\ref{sec:fairness}.}
\label{fig:pagination}
\end{figure}

\begin{lstlisting}[language=Python, caption={The stop condition, the part worth reading twice}]
            subs = data.get('result', [])
            if not subs:
                break
            all_subs.extend(subs)
            time.sleep(delay)
            if len(subs) < batch_size:
                break
            start += batch_size
\end{lstlisting}

There are two stop conditions, and they are not redundant:

\begin{enumerate}
\item \textbf{An empty page} means we asked past the end. This catches the exact-multiple
case, where a user has precisely 20{,}000 submissions and the second page is full.
\item \textbf{A short page} means this page \emph{is} the end. This is the common case, and
it saves one wasted round trip that would have returned nothing.
\end{enumerate}

\begin{keyidea}{Why ``short page means done'' is safe here}
This inference is not universally valid. It is valid against an API that only returns a
short page when it has run out of records, and never returns a short page for its own
internal reasons (throttling, partial results). Codeforces behaves that way, so the
optimisation holds. I confirmed the branch is real and load-bearing by measuring live:
\code{tourist} returned 5459 records on page one, which is shorter than the 10{,}000
requested, so the loop stopped after a single request instead of making a second one.
\end{keyidea}

\subsection{\code{fetch\_problemset}: one request, the whole catalogue}

The counterpart is trivial by comparison: a single request for the entire problem
catalogue, no pagination, a 60 second timeout rather than 30 because the response is much
larger. On success it sleeps and returns the problems list; on any failure it prints and
returns an empty list.

\begin{honest}{Failing to fetch the catalogue produces a successful-looking run}
If this request fails, the function returns \code{[]}. That empty list flows through
\code{get\_filtered\_problems}, which happily filters nothing, into
\code{save\_to\_excel}, which prints ``No problems to save.'' and returns. Then
\code{\_\_main\_\_} prints ``Done.'' and the process exits with status 0, the code meaning
success. A scheduled job wrapping this script would see a green tick and no spreadsheet.
The word ``Done.'' printing after a total failure to fetch anything is the clearest
example in the codebase of the exit-code problem discussed in section~\ref{sec:honest}.
\end{honest}

\subsection{\code{collect\_solved}: two callers, one rule}

\begin{lstlisting}[language=Python, caption={The normalisation step, in full}]
def collect_solved(handles, delay):
    """Set of (contestId, index) pairs with an OK verdict across the given handles."""
    solved = set()
    for u in handles:
        for s in fetch_user_status(u, delay):
            if s.get('verdict') == 'OK':
                prob = s.get('problem', {})
                cid = prob.get('contestId') or s.get('contestId')
                idx = prob.get('index')
                if cid and idx:
                    solved.add((cid, idx))
    return solved
\end{lstlisting}

Nine lines, and three separate decisions inside them:

\begin{enumerate}
\item \textbf{\code{verdict == 'OK'}} is the fairness rule itself. A member who attempted a
problem and failed has not solved it, and the problem stays eligible. Every other verdict
falls through.
\item \textbf{\code{prob.get('contestId') or s.get('contestId')}} is defensive
normalisation. The contest id normally lives on the nested problem, but not always, so the
code falls back to the submission envelope. This is a real quirk of the API surface, not
speculative hardening.
\item \textbf{\code{if cid and idx}} refuses to build a half-formed key. Without it, an
entry missing an index would insert \code{(2236, None)} into the set, a key that can never
match a real problem but silently pollutes the data.
\end{enumerate}

\begin{keyidea}{Why this function exists at all}
Include and exclude need the identical computation over different handle lists. The
README records that this was originally one loop copy-pasted twice, so a fix to the
verdict rule had to be made in two places or it was wrong in one. Parameterising by
\code{handles} and calling it twice means the \code{OK} rule, the fallback, and the guard
each exist exactly once. This is the clearest refactor in the project, and it is worth
naming as such: the duplication was not merely ugly, it was a correctness hazard.
\end{keyidea}

\begin{honest}{A \code{falsy} guard where a \code{None} check was meant}
\code{if cid and idx} rejects any value Python considers falsy, not just missing ones. A
contest id of \code{0} would be silently discarded. No Codeforces contest has id 0, so
this is harmless in practice and I am flagging it as a code-reading observation rather
than a defect: the guard expresses ``is present'' using a test that means ``is truthy'',
and those two only coincide because of a fact about the data that is nowhere written down.
\end{honest}

\subsection{\code{get\_filtered\_problems}: bucketing, then the three gates}

This orchestrates everything. First it buckets the catalogue by rating:

\begin{lstlisting}[language=Python]
    by_rating = {r: [] for r in args.ratings}
    for p in problems:
        r = p.get('rating')
        if r in by_rating:
            by_rating[r].append(p)
\end{lstlisting}

\begin{keyidea}{Pre-seeding the dictionary is doing two jobs}
\code{by\_rating} is created with a key for every requested rating and an empty list
behind it, \emph{before} any problem is examined. That means the subsequent
\code{if r in by\_rating} test is simultaneously ``do we want this rating?'' and ``where
does it go?''. One dictionary lookup replaces both a membership test against the requested
ratings and a create-if-absent. It also guarantees a requested rating with zero matching
problems still appears downstream as an empty bucket rather than vanishing.
\end{keyidea}

Problems whose rating is unset are dropped here automatically, because \code{p.get('rating')}
returns \code{None} and \code{None} is not a key in \code{by\_rating}. That case is real: I
measured the live catalogue at 11{,}295 problems, of which \textbf{289 carry no rating at
all}. The \code{.get} is doing genuine work.

Then the two solved-sets are built, each guarded so that disabling a filter also skips its
download entirely, and each problem faces three gates:

\begin{lstlisting}[language=Python, caption={The core predicate}]
            key = (p['contestId'], p['index'])
            if (
                (args.no_include_filter or key in include_solved)
                and (args.no_exclude_filter or key not in exclude_solved)
                and passes_tag_filter(p)
                ):
                valid.append(p)
\end{lstlisting}

The \code{or} short-circuits are how a disabled filter is expressed: if
\code{no\_include\_filter} is true, the whole clause is true without ever consulting the
(empty) set. Correct, and it reads cleanly.

\begin{honest}{Inconsistent defensiveness on the same two fields}
\code{collect\_solved} carefully uses \code{.get} with a fallback for \code{contestId} and
guards against missing values. Twenty lines later, \code{get\_filtered\_problems} uses
direct subscripting, \code{p['contestId']} and \code{p['index']}, which raises
\code{KeyError} and kills the run if either is absent. The same two fields are treated as
untrustworthy in one function and guaranteed in the other.

I checked whether this can actually fire: across all 11{,}295 problems in the live
catalogue, \textbf{zero} are missing \code{contestId} and \textbf{zero} are missing
\code{index}. So the crash is unreachable today. The defect is the inconsistency itself,
not a live bug: a reader cannot tell from the code which function is right about the API's
guarantees, and the answer turns out to be that the defensive one is defending against a
quirk of \emph{submissions}, not of the catalogue. That is a reasonable distinction which
the code never states.
\end{honest}

\subsection{\code{passes\_tag\_filter}: four modes, one asymmetry}

Defined as a closure inside \code{get\_filtered\_problems} so it can read \code{args}
without being passed it.

\begin{figure}[H]
\centering
\resizebox{0.9\textwidth}{!}{
\begin{tikzpicture}[node distance=8mm]
  \node[box] (in) {problem's tag set $T$};
  \node[dec, below=of in] (exc) {$T \cap$ exclude-tags\\non-empty?};
  \node[box, right=24mm of exc, draw=warnred] (rej1) {\textbf{reject}};
  \node[dec, below=13mm of exc] (has) {include-tags given\\and mode $\neq$ \code{off}?};
  \node[box, left=20mm of has, fill=softgrey] (acc1) {\textbf{accept}};
  \node[dec, below=13mm of has] (mode) {which mode?};
  \node[box, below left=12mm and 16mm of mode] (some) {\code{some}\\$T \cap I \neq \emptyset$};
  \node[box, below=12mm of mode] (all) {\code{all}\\$I \subseteq T$};
  \node[box, below right=12mm and 16mm of mode] (strict) {\code{strict}\\$T = I$};
  \node[box, below=10mm of all, fill=softgrey] (acc2) {\textbf{accept} if satisfied,\\else \textbf{reject}};

  \draw[flow] (in) -- (exc);
  \draw[flow] (exc) -- node[lbl] {yes} (rej1);
  \draw[flow] (exc) -- node[lbl] {no} (has);
  \draw[flow] (has) -- node[lbl] {no} (acc1);
  \draw[flow] (has) -- node[lbl] {yes} (mode);
  \draw[flow] (mode) -- (some);
  \draw[flow] (mode) -- (all);
  \draw[flow] (mode) -- (strict);
  \draw[flow] (some) |- (acc2);
  \draw[flow] (all) -- (acc2);
  \draw[flow] (strict) |- (acc2);
\end{tikzpicture}}
\caption{Tag filtering. $I$ is the include-tags set. Note the asymmetry at the top: the
exclude-tags gate runs unconditionally, before and independent of the mode.}
\end{figure}

\begin{keyidea}{The asymmetry is deliberate and correct}
\code{--tag-mode off} disables the \emph{include} tags but leaves \code{--exclude-tags}
fully active. That is the right semantics. ``Mode'' answers the question ``how should I
match the tags you asked for?'', which is meaningless for tags you asked to avoid: there
is only one way to avoid a tag. Reading the code, the exclude check sits above the mode
logic entirely, so the two can never interfere. Someone could reasonably expect
\code{off} to mean ``no tag filtering at all'', and it does not; the README documents the
behaviour correctly.
\end{keyidea}

The three include modes, precisely:
\begin{itemize}
\item \code{some}: at least one requested tag present. The loosest. \emph{Give me
anything greedy-ish.}
\item \code{all}: every requested tag present, extras allowed. \emph{Must be greedy AND
maths, whatever else it is.}
\item \code{strict}: the problem's tag set equals the requested set exactly, no extras.
\emph{Purely greedy, nothing else.} On real data this is extremely restrictive, since most
problems carry three or more tags.
\end{itemize}

\subsection{\code{pick\_choices\_per\_rating}: two lines, one real weakness}

\begin{lstlisting}[language=Python]
def pick_choices_per_rating(filtered, n):
    selection = {}
    for r, plist in filtered.items():
        selection[r] = plist[:n]
    return selection
\end{lstlisting}

\begin{honest}{``Selection'' is a truncation, and it is biased, not neutral}
\code{plist[:n]} takes the first N in whatever order the API returned. It does not sample.
Codeforces returns the problemset newest-first, so this systematically hands you the most
recent problems and never surfaces an older classic, and two runs with the same filters
produce nearly the same sheet, which undercuts the point of a tool for a recurring study
group.

This is not theory. In my live run (section~\ref{sec:verified}) I asked for three problems
each at ratings 800, 1200 and 1600. The contest ids that came back were 2237, 2236, 2236,
2227, 2189, 2178, 2207, 2172, 2157: all from a narrow, recent band, out of a catalogue
spanning contest ids from the low hundreds upward. The bias is total, not marginal.

The README's ``What I'd do differently'' names this and proposes \code{random.sample} with
a seed. That is the correct fix and it is about four lines. Worth being ready to say in an
interview why it was not done: the honest answer is that the tool met the group's need as
written, not that the problem was unrecognised.
\end{honest}

\subsection{\code{save\_to\_excel}: flatten and write}

Takes the nested dictionary of rating to problem-list and flattens it to one row per
problem, deriving a clickable link from the identity tuple. It bails out early with ``No
problems to save.'' rather than writing an empty file.

\begin{jargon}{pandas and DataFrame}
\code{pandas} is Python's standard library for tabular data. A \code{DataFrame} is a table
in memory with named columns, essentially a spreadsheet as a variable. Here it is used at
its most basic: build a list of dictionaries where each dictionary is a row, hand it to
\code{pd.DataFrame}, and call \code{.to\_excel}. The heavy lifting of writing the actual
\file{.xlsx} binary format is delegated to \code{openpyxl}.
\end{jargon}

\begin{honest}{pandas is a heavy dependency for what it does here}
The only pandas features used are ``list of dicts becomes a table'' and ``table becomes
\file{.xlsx}''. \code{openpyxl} alone can do that, and pandas depends on \code{numpy}, so
this pulls tens of megabytes of scientific computing to write a nine-row sheet. It is a
defensible convenience choice, and \code{to\_excel} genuinely is one line, but it is worth
being able to name the tradeoff rather than being caught by it.
\end{honest}

\section{Dependencies}

\begin{table}[H]
\centering\small
\begin{tabular}{@{}llp{7cm}@{}}
\toprule
\textbf{Package} & \textbf{Pin} & \textbf{Role} \\
\midrule
\code{requests} & \code{>=2.31,<3} & Every HTTP call to the Codeforces API. Used for its
timeout support and \code{raise\_for\_status}. \\
\code{pandas} & \code{>=2.0,<3} & Builds the table and writes the sheet. \\
\code{openpyxl} & \code{>=3.1,<4} & The actual \file{.xlsx} writer. Never imported
directly; pandas requires it to be installed to write Excel. \\
\bottomrule
\end{tabular}
\caption{All three dependencies. Verified installable and working on Python 3.12.3.}
\end{table}

\begin{keyidea}{Why \code{openpyxl} appears in requirements but never in the code}
A beginner grepping for \code{import openpyxl} will not find it and may conclude the
dependency is spurious. It is not. \code{pandas.to\_excel} is a thin front end that looks
up an engine at runtime and delegates; without \code{openpyxl} installed it raises at the
moment of writing, after all the API downloading is already done. Pinning it explicitly is
correct: it is a real runtime requirement that the import graph does not reveal.
\end{keyidea}

The pins are compatible-range style: at least a known-good minor, below the next major.
This is the sensible middle ground between unpinned (breaks silently when an upstream major
lands) and exact pins (rots, needs constant bumping).

\section{Verified behaviour}
\label{sec:verified}

I did not take the README's claims on trust. I built a clean virtual environment on Python
3.12.3, installed the three pinned dependencies, and ran the exact command the README
documents, against the live Codeforces API.

\begin{lstlisting}[language=bash, caption={The command from the README, run verbatim}]
python generate_mashup.py \
    --include-users ksun48 \
    --exclude-users salmanadnan2025 Musab1Blaser \
    --ratings 800 1200 1600 \
    --choices-per-rating 3 \
    --include-tags greedy --tag-mode some \
    -o mashup_test.xlsx
\end{lstlisting}

\begin{table}[H]
\centering\small
\begin{tabular}{@{}p{5.6cm}ll@{}}
\toprule
\textbf{README claim} & \textbf{Measured} & \textbf{Verdict} \\
\midrule
Takes about 14 seconds & 12.1 seconds & \textcolor{okgreen}{\textbf{holds}} \\
Produced 9 rows & 9 rows & \textcolor{okgreen}{\textbf{holds}} \\
3 per rating & 3 at each of 800, 1200, 1600 & \textcolor{okgreen}{\textbf{holds}} \\
All tagged \code{greedy} & 9 of 9 carry \code{greedy} & \textcolor{okgreen}{\textbf{holds}} \\
46 trusted handles, 34 group members & 46 and 34, no duplicates & \textcolor{okgreen}{\textbf{holds}} \\
Bad handles print and skip, sheet still generates & reproduced with a fabricated handle & \textcolor{okgreen}{\textbf{holds}} \\
\bottomrule
\end{tabular}
\caption{Every checkable README claim, checked. This README is unusually accurate; nothing
in it was found to be overstated.}
\end{table}

The output columns are \code{rating}, \code{contestId}, \code{index}, \code{name},
\code{tags}, \code{link}, exactly as documented. A representative row: rating 800, contest
2236, index A, ``Games on the Train'', tags \code{greedy, math}, with a working problemset
link.

I also confirmed the include and exclude default lists are disjoint (no handle appears in
both), so the pathological case of a handle being downloaded twice and simultaneously
qualifying and disqualifying every problem it has solved does not arise on a default run.
Nothing in the code prevents it if a user constructs it by hand.

\section{Performance, measured}
\label{sec:perf}

The README says a default run pulls ``hundreds of megabytes''. I checked rather than
repeating it. Fetching the first page of submissions for three of the default trusted
handles:

\begin{table}[H]
\centering\small
\begin{tabular}{@{}lrr@{}}
\toprule
\textbf{Handle} & \textbf{Records, page 1} & \textbf{Wire size} \\
\midrule
\code{tourist}  & 5{,}459  & 3.1 MB \\
\code{ksun48}   & 9{,}101  & 5.2 MB \\
\code{jiangly}  & 10{,}000 (page 1 of more) & 5.6 MB \\
\bottomrule
\end{tabular}
\caption{Measured live. \code{jiangly} hit the 10{,}000 cap, so that 5.6 MB is a floor,
not a total.}
\end{table}

Averaging those three and extrapolating across the 46 default include handles gives
roughly \textbf{214 MB, as a lower bound}, since it counts only one page each and ignores
the 34 exclude handles entirely. The README's ``hundreds of megabytes'' is
\textbf{substantiated, and if anything conservative}. All of it is downloaded to answer a
set-membership question, and none of it is cached: run the tool twice in a row with
identical arguments and it downloads everything again.

\begin{jargon}{Caching}
Storing the result of expensive work so a later run can reuse it instead of redoing it.
Here, writing each handle's solved-set to a small file on disk would turn the second run
of a day from a 200 MB download into reading a few kilobytes. The tool does none of this,
which is fine for a script run occasionally by one person and would be the first thing to
fix if it ran more often.
\end{jargon}

\section{Honest assessment}
\label{sec:honest}

The code is clean, readable, and does what it claims. The findings below are real, and
several are already named in the README's own ``What I'd do differently'', which is a
point in the author's favour: he found most of them himself.

\subsection{The one that actually matters}
\label{sec:fairness}

\begin{honest}{A failed fetch silently voids the fairness guarantee}
This is the most serious defect in the project, because it defeats the tool's entire
purpose while looking like it worked.

Follow figure~\ref{fig:pagination}: every error path in \code{fetch\_user\_status} prints
a message and returns whatever it has, which for a failure on the first request is an
empty list. \code{collect\_solved} cannot distinguish ``this member has solved nothing''
from ``we failed to find out what this member has solved''. Both are the empty set.

So if a network blip, a rate limit, or a renamed handle causes one group member's fetch to
fail, that member's solved problems are treated as unsolved, they sail through the exclude
filter, and the member is handed a problem they have already done. That is precisely the
scenario the tool exists to prevent. The run prints one line among many and exits 0.

I reproduced the mechanism live: an invented handle in \code{--exclude-users} produced
\code{Error fetching submissions for ...: 400 Client Error}, and then the sheet generated
normally and the process exited 0. The README notes that this happened for real, with two
stale handles in the default exclude list returning 400.

The severity is asymmetric and the code does not reflect that. A failed \emph{include}
fetch only costs you some candidate problems, which is a quality loss. A failed
\emph{exclude} fetch costs you the guarantee, which is a correctness loss. They are handled
identically. A failed exclude fetch should abort the run. The README's own ``What I'd do
differently'' reaches the same conclusion, in the same words: fairness is the whole point.
\end{honest}

\subsection{The rest, in descending order of consequence}

\begin{honest}{No tests, at all}
There is no test file in the repository and no test dependency. Every claim about this
code, including every one in this document, had to be established by reading it or by
running it against the live API. The pure functions are the easy, obvious targets:
\code{passes\_tag\_filter}'s four modes and \code{pick\_choices\_per\_rating} are
completely deterministic, take plain data, touch no network, and could be pinned down in a
few dozen lines. \code{collect\_solved}'s envelope-normalisation is exactly the kind of
fiddly logic that regresses silently. As it stands, the only way to know a refactor was
safe is to run the whole thing and eyeball a spreadsheet.
\end{honest}

\begin{honest}{Everything exits 0, and reports through \code{print}}
There is no logging and no exit code discipline. Success, partial data, a completely failed
catalogue fetch, and a run that saved nothing all exit 0 and all report through
\code{print} to stdout. Errors are interleaved with results on the same stream. The script
is therefore not safely automatable: nothing downstream can detect that it failed. Adding
\code{logging} to stderr and a nonzero exit when no rows were written is small work with a
large payoff, and the README names it.
\end{honest}

\begin{honest}{Dead data at the project root}
\code{all\_users.txt}, covered in full earlier: unread by any code, and demonstrably a
stale snapshot of the script's own constants, six handles behind on the exclude list. It
should be deleted.
\end{honest}

\begin{honest}{The default run is close to unusable, and the defaults advertise it}
A bare \code{python generate\_mashup.py} downloads the full submission history of 80
handles, on the order of a quarter of a gigabyte at minimum (section~\ref{sec:perf}), with
a 0.5 second sleep between every page, to produce a spreadsheet. The README warns that
defaults ``take a while''. That warning is doing a lot of work: this is measured in many
minutes, and the tool's own documentation only ever demonstrates it on a cut-down example.
The defaults reproduce a specific study group's configuration, which is defensible as
provenance, but it means the out-of-box experience is the slowest possible path.
\end{honest}

\begin{honest}{A personal handle is baked into a public default}
\code{salmanadnan2025} sits in \code{DEFAULT\_EXCLUDE\_USERS} in a public repository.
That is intentional and harmless (he was in the study group, and it is a public handle on a
public site), but it is worth being conscious that the shipped defaults are one particular
group's roster, including the author's own handle, rather than anything general. Anyone
cloning this runs a filter built around 34 strangers unless they pass their own flags.
\end{honest}

\begin{honest}{The \code{--delay} flag guards the wrong boundary}
\code{--delay} sleeps between pages within a handle and after the problemset fetch, but
Codeforces rate limits per client, not per handle. The sleep placement is a reasonable
approximation and it evidently works, but it is not derived from the published limit. If a
run ever does get throttled, the failure mode is the one in
section~\ref{sec:fairness}: a non-\code{OK} response, a printed line, and a silently
weakened exclude set.
\end{honest}

\section{What to say about this project}

\begin{keyidea}{The defensible core}
Strip away the polish and there is one genuinely good idea here, and it generalises well
beyond Codeforces: \textbf{when an API hands you the same entity in two different envelope
shapes, do not try to compare the objects. Choose a canonical identity, reduce both sides
to it, and let set arithmetic do the rest.} That is what makes the include and exclude
filters line up, it is why there is no nested search in the code, and it is the thing worth
leading with.

The second defensible point is the \code{collect\_solved} extraction: recognising that two
copy-pasted loops were not merely repetitive but a correctness hazard, since the verdict
rule could drift between them, and collapsing them into one parameterised function.

The honest framing of the whole project: it is a small, well-written tool that solved a
real problem for a real group, with one significant unfixed flaw (a failed exclude fetch
silently voids the guarantee), one significant inefficiency (downloading hundreds of
megabytes to answer a membership question that \code{solvedCount} would answer for free),
and no tests. The author identified most of that himself before anyone asked, which is the
best thing in the repository.
\end{keyidea}

\end{document}
